\subsection{Statement}
\label{jep104local035}

We consider the functor
\[
\Phi : D^{\mathrm{b}}_{G_{\jepPubScript{O}}}(\mathrm{Gr}, \Bbbk) \to D^{\mathrm{b}}_{\jepPubEuler{IW}}(\mathrm{Gr}, \Bbbk)
\]
defined by
\[
\Phi(\jepPubEuler{F}) = \Delta^{\jepPubEuler{IW}}_\varsigma(\Bbbk) \star^{G_{\jepPubScript{O}}} \jepPubEuler{F}.
\]
In view of~\eqref{jep104local003} (or, alternatively, arguing as in~\cite{jep104ref17}), in this definition $
\Delta^{\jepPubEuler{IW}}_{\varsigma}(\Bbbk)$ can be replaced by $\nabla^{\jepPubEuler{IW}}_\varsigma(\Bbbk)$ or $ \mathrm{IC}^{\jepPubEuler{IW}}_\varsigma(\Bbbk)$; in particular this shows that the conjugate of $\Phi$ by Verdier duality is the similar functor using the character $\psi^{-1}$ instead of $\psi$. 

\begin{lem}
\label{jep104local016}
 The functor $\Phi$ is t-exact for the perverse t-structures.
\end{lem}

\begin{proof}
In the case where $\Bbbk$ is a field, the claim follows from Lemma~\ref{jep104local018}. The general case follows using an $\ell$-modular triple $(\jepPubBbb{K},\jepPubBbb{O},\jepPubBbb{L})$ as in Section~3.2 with $\Bbbk=\jepPubBbb{O}$, and the associated extension of scalars functors. Namely, if $\jepPubEuler{F}$ is in $\mathrm{Perv}_{G_{\jepPubScript{O}}}(\mathrm{Gr},\jepPubBbb{O})$, then
\[
\jepPubBbb{K}\jepDerivedTensor_{\jepPubBbb{O}}\Phi(\jepPubEuler{F})\cong\Phi(\jepPubBbb{K}\jepDerivedTensor_{\jepPubBbb{O}}\jepPubEuler{F})
\]
is perverse; hence any perverse cohomology object ${}^{p}\jepPubEuler{H}^i(\Phi(\jepPubEuler{F}))$ with $i\neq0$ is torsion. On the other hand,
\[
\jepPubBbb{L}\jepDerivedTensor_{\jepPubBbb{O}}\Phi(\jepPubEuler{F})\cong\Phi(\jepPubBbb{L}\jepDerivedTensor_{\jepPubBbb{O}}\jepPubEuler{F})
\]
lives in perverse degrees $-1$ and $0$ since $\jepPubBbb{L}\jepDerivedTensor_{\jepPubBbb{O}}\jepPubEuler{F}$ lives in these degrees. If ${}^{p}\jepPubEuler{H}^i(\Phi(\jepPubEuler{F}))$ were nonzero for some $i>0$, then taking $i$ maximal with this property we would obtain that ${}^{p}\jepPubEuler{H}^i(\jepPubBbb{L}\jepDerivedTensor_{\jepPubBbb{O}}\Phi(\jepPubEuler{F}))\neq0$, a contradiction. On the other hand, if ${}^{p}\jepPubEuler{H}^i(\Phi(\jepPubEuler{F}))$ was nonzero for some $i<0$, then taking $i$ minimal with this property we would obtain that ${}^{p}\jepPubEuler{H}^{i-1}(\jepPubBbb{L}\jepDerivedTensor_{\jepPubBbb{O}}\Phi(\jepPubEuler{F}))\neq0$, a contradiction again.
\end{proof}
 

We will denote by
\[
 \Phi^0 : \mathrm{Perv}_{G_{\jepPubScript{O}}}(\mathrm{Gr}, \Bbbk) \to \mathrm{Perv}_{\jepPubEuler{IW}}(\mathrm{Gr}, \Bbbk)
\]
the restriction of $\Phi$ to the hearts of the perverse t-structures, so that $\Phi^0$ is an exact functor between abelian categories.
The main result of this section is the following theorem, whose proof will be given in the next subsection.

\begin{thm}
\label{jep104local038}
The functor $\Phi^0$ is an equivalence of categories.
\end{thm}

